% §7 Open problems
Five problems stay open, in increasing order of ambition.

\paragraph{Learned emission mapping.}
The rule-based classifier is the interface's seed, not its ceiling.
A learned state estimator --- trained on pattern-table labels, then corrected --- would
replace \texttt{infer\_states()} with no downstream change; the open question is what its
uncertainty does to coverage (a walk with a soft state distribution induces a coverage
distribution, not a number).

\paragraph{Calibrated embeddings as features.}
The embedding variant of Section~\ref{sec:evaluation} swaps the feature step; the stronger
version conditions the projection itself on calibrated user embeddings \citep{cue2026}, so
the behavior plane separates conversations by calibrated semantic distance rather than
structural counts.

\paragraph{Coverage beyond two dimensions.}
The 2-D projection is a visualization-driven commitment; the coverage \emph{definition} does
not need it.
Kernel two-sample tests over the full feature space (MMD-style) would measure region coverage
without any projection, at the cost of the atlas and its exemplar loop --- the open problem
is a formulation that keeps per-region exemplars in high dimensions.

\paragraph{Closed-loop simulator improvement.}
The recommendation loop is currently human-in-the-middle: uncovered exemplars become scenario
seeds by an analyst's decision.
Closing the loop --- uncovered region $\to$ generated scenario $\to$ coverage delta, with no
human step --- is the simulator-building agenda \citep{cue2026,distributionalGap2026} run on
the atlas as its objective function.

\paragraph{Voice-physics channels.}
Barge-in, endpointing behavior, and tool round-trip latency die at the transport boundary of
current observability stacks.
The companion grammar specifies a \texttt{voice.*} semantic-convention extension
\citep{garcia2026grammar}; the atlas extends to it unchanged (the channels become feature
dimensions), making the missing piece a transport-side export fix rather than a method
problem.
